\documentclass{article}

\usepackage[utf8]{inputenc}     % pdfLaTeX
\usepackage[T1]{fontenc}
\usepackage[magyar]{babel}
\usepackage{amsmath}

% Set page size and margins
% Replace `letterpaper' with `a4paper' for UK/EU standard size
\usepackage[a4paper,top=2cm,bottom=2cm,left=3cm,right=3cm,marginparwidth=1.75cm]{geometry}
%for arrows
\usepackage{textcomp}

\title{Személynevek helyesírása}
\author{Lengyel Zoltán Péter\thanks{Csukás Ibolya tanítása alapján}}
\date{Legutóbb frissítve: 2026. Március 15.}

\begin{document}

\maketitle
\tableofcontents
\newpage

\begin{itemize}
    \item Egyes családnevek tartalmazhatnak régies magánhangzót és mássalhangzót
    \begin{itemize}
        \item pl. Csoóri; Weöres; Damjanich; Batthyány
        \item Itt figyelni kell az elválasztásra, a régies betűket nem választhatjuk el!
        \begin{itemize}
            \item Csoó-ri, Weö-res, Bat-thyá-ny
        \end{itemize} 
    \end{itemize}
    \item A toldalékot általában egybeírjuk a személynévvel
    \begin{itemize}
        \item pl. Vörösmartyt; Damjanichot
    \end{itemize}
    \item Régies betűre végződő névhez a -val / -vel rag kiejtés szerint hasonul
    \begin{itemize}
        \item pl. Kossuthtal; Babitscsal
    \end{itemize}
    \item Hosszú mássalhangzóra végződő neveknél a -val / -vel hasonult ragot kötőjellel írjuk
    \begin{itemize}
        \item pl. Kiss-sel; Wass-sal; Ivett-tel; Marcell-lel
        \item Ez más toldalékokkal is így van
        \begin{itemize}
            \item Mariann-nak; Ivett-től
        \end{itemize}
    \end{itemize}
    \item Néma hangra végződő nevekhez a toldalékot a kiejtés szerint, kötőjellel kapcsoljuk
    \begin{itemize}
        \item pl. Shakespeare-rel; Anne-ről
    \end{itemize}

    \item -i képzős alakok: 
    \begin{itemize}
        \item A teljes névhez kötőjellel kapcsoljuk, nagybetűk megmaradnak
        \begin{itemize}
            \item pl. Petőfi Sándor-os
        \end{itemize}
        \item Csak a vezeték/családnévhez kötőjel nélkül kapcsoljuk, kisbetűvel írjuk
        \begin{itemize}
            \item pl. petőfies; kossuthi
        \end{itemize}
    \end{itemize}
\end{itemize}


\end{document}